\chapter[p-adic Fields]{$p$-adic Fields}

\section[p-adic Equations]{$p$-adic Equations}\label{sec:padic_equations}

\begin{definition}\label{def:primitive}
\lean{Function.IsPrimitive}
\leanok
Let $R$ be a ring. An element $(x_1, \cdots, x_n) \in R^n$ is \emph{primitive} if at least one of the $x_i$ is invertible.
\end{definition}


\begin{proposition}\label{prop:common_root}
\lean{Function.IsPrimitive}
\leanok
\uses{def:primitive}
Let $f^{(i)} \in \Z_p[X_1, \cdots, X_m]$ be homogeneous polynomials with $p$-adic integer coefficients. For each $n \ge 1$, denote by $f_n^{(i)}$ the reduction of $f^{(i)}$ modulo $p^n$.
 The following are equivalent:
\begin{enumerate}[a)]
\lean{Padic.common_root_tfae}
\leanok
\item The $f^{(i)}$ have a nontrivial common zero in $\Q_p^m$.
\item The $f^{(i)}$ have a common primitive zero in $\Z_p^m$.
\item For all $n \ge 1$, the $f_n^{(i)}$ have a common primitive zero in $(\Z / p^n \Z_p)^m$.
\end{enumerate}
\end{proposition}

\begin{proof}
The implication (b) $ \implies$ (a) is trivial. Conversely, if $x = (x_1, \cdots, x_m) \in \Q_p^m$ is a nontrivial common zero of the $f^{(i)}$, put
\[ h := \min(v_p(x_1), \cdots, v_p(x_m)) \quad \text{and} \quad y := p^{-h} x. \]
Then $y \in \Z_p^m$ is a common primitive zero of the $f^{(i)}$.
For the equivalence (b) $\iff$ (c), Serre applies the first Lemma from \cite[\S II.2.1]{Ser73}. However, using \mathlib's \texttt{DirectLimit} definition here might be tricky. It could be easier to just prove this specific result; note that the file \texttt{Mathlib.NumberTheory.Padics.RingHoms} contains useful results relating $\Z_p$ and $\Z/p^n\Z$ (without relying on \texttt{DirectLimit}).

\end{proof}

\section[Squares in the p-adic units]{Squares in $\Q_p^{\times}$}\label{sec:padic_squares}

\begin{theorem}\label{thm:padic_squares_open}
\lean{Padic.unitSquares}
\leanok
${\Q_p^\times}^2$ is an open subgroup of $\Q_p^\times$.
\end{theorem}

\begin{proof}
\leanok
The only nontrivial part is to show that the squares form an open set.  Since $\Q_p$ is a field, the unit topology agrees with the subspace topology, so it suffices to check that the set of nonzero squares in $\Q_p$ is open. Define $r$ as $1$ if $p$ is odd, or as $2^{-2}$ if $p = 2$, and denote by $B_{1, r}$ the open ball with center $1$ and radius $r$. For each $c \in B_{1, r}$, by applying Hensel's lemma with the polynomial $X^2 - c$, we obtain that $c$ is a square.
Then, for each nonzero square $s \in \Q_p$, $s \cdot B_{1, r}$ is an open neighbourhood of $s$ contained in the set of nonzero squares, which shows that the latter is open.
See \cite[\S II.3.3]{Ser73} for more details.
\end{proof}

\begin{lemma}
  \label{lem:multivariable_hensel}
  \lean{PadicInt.multivariable_hensel}
  \leanok
Let $p$ be a prime, let $f \in \Z_p[X_1, \cdots, X_m]$ be a polynomial with $p$-adic integer
coefficients, and let $x \in (\Z/p\Z)^m$ be a zero of $f$ modulo $p^n$, for some positive integer
$n$. Suppose that, for some $j \in \{1, \cdots, m\}$, the partial derivative
$\frac{\partial f}{\partial X_j}(x)$ is nonzero modulo $p^k$, where $0 < 2k < n$.
Then there exists a zero $z$ of $f$ in $\Z_p^m$ whose reduction modulo $p^n$ is $x$, and such that
$z-x$ is divisible by $p^{n-k}$.
\end{lemma}

\begin{proof}
This is in Serre's \cite[\S II.2 Theorem 1]{Ser73}.
\end{proof}

\section[Units and Legendre symbols]{Units and Legendre symbols}\label{sec:padic_legendre}

\begin{lemma}\label{lem:padic_unit_part_unit}
\lean{Padic.norm_mul_pow_neg_valuation_eq_one}
\leanok
Let $p$ be a prime, let $x \in \Q_p^\times$, and let $v_p$ denote the $p$-adic valuation. Then $\|x \cdot p^{-v_p(x)}\| = 1$.
\end{lemma}
\begin{proof}
\leanok
Straightforward.
\end{proof}

\begin{definition}\label{def:padicUnitPart}
  \lean{Padic.unitPart}
  \uses{lem:padic_unit_part_unit}
  \leanok
Let $p$ be a prime and $x \in \Q_p^\times$. The \emph{unit part} of $x$ is defined as the element
of $\Z_p^\times$ corresponding to $x \cdot p^{-v_p(x)}$.
\end{definition}

\begin{definition}\label{def:legendreSym}
  \lean{PadicInt.legendreSym}
  \leanok
Let $p$ be a prime and $u \in \Z_p^\times$. The \emph{Legendre symbol} $\left(\dfrac{u}{p}\right)$
is defined as
\[\left(\dfrac{u}{p}\right)=\left(\dfrac{\overline{u}}{p}\right),\]
where $\overline{u}$ is the image of $u$ under the natural map $\Z_p^\times \to (\Z/p\Z)^\times$.
\end{definition}
